% !TEX program = lualatex
\documentclass[professionalfonts]{beamer}
\newcommand{\slidemode}{1}  % カラー投影モード
\usepackage{beamer_template}
\usepackage{enumerate}

\newcommand{\LectureNum}{13}
\newcommand{\LectureTitle}{フーリエ変換の性質と公式}
\newcommand{\CourseName}{応用数学}
\renewcommand{\TermName}{前期}

\begin{document}

\begin{frame}{第13回　フーリエ変換の性質と公式}
  \begin{block}{今回の内容}
  フーリエ変換の性質と公式
  \end{block}
  \vfill\hfill{\scriptsize \textcolor{gray}{第3章 フーリエ解析　§2}}
\end{frame}

\begin{frame}{フーリエ変換の性質}
  \begin{block}{}
  \begin{tabular}{ll}
    I.   & 線形性：$\mathcal{F}[c_{1}f_{1}+c_{2}f_{2}] = c_{1}F_{1}(u)+c_{2}F_{2}(u)$\\[4pt]
    II.  & 移動則（時間）：$\mathcal{F}[f(x-a)] = e^{-iau}F(u)$\\[4pt]
    III. & 移動則（周波数）：$\mathcal{F}[e^{iax}f(x)] = F(u-a)$\\[4pt]
    IV.  & 相似性：$\mathcal{F}[f(ax)] = \dfrac{1}{|a|}F\!\left(\dfrac{u}{a}\right)$\\[6pt]
    V.   & 微分則：$\mathcal{F}[f^{(n)}(x)] = (iu)^{n}F(u)$\\[4pt]
    VI.  & 双対微分則：$\mathcal{F}[(-ix)^{n}f(x)] = F^{(n)}(u)$\\
  \end{tabular}
  \end{block}
\end{frame}

\begin{frame}{性質の証明（II・IV・V）}
  $II:\; \mathcal{F}[f(x-a)] = \displaystyle\int_{-\infty}^{\infty} f(x-a)e^{-iux}\,dx$\\[2pt]
  $t=x-a$ と置換：$= \displaystyle\int_{-\infty}^{\infty} f(t)e^{-iu(t+a)}\,dt = e^{-iau}\int_{-\infty}^{\infty}f(t)e^{-iut}\,dt = e^{-iau}F(u)$\\[8pt]
  $IV:\; \mathcal{F}[f(ax)] = \displaystyle\int_{-\infty}^{\infty} f(ax)e^{-iux}\,dx$\\[2pt]
  $t=ax$ と置換（$a>0$）：$= \dfrac{1}{a}\displaystyle\int_{-\infty}^{\infty}f(t)e^{-i(u/a)t}\,dt = \dfrac{1}{a}F\!\left(\dfrac{u}{a}\right)$\\[4pt]
  （$a<0$ のとき符号が逆になり同様に $\dfrac{1}{|a|}F\!\left(\dfrac{u}{a}\right)$）\\[8pt]
  $V:\; \mathcal{F}[f'(x)] = \displaystyle\int_{-\infty}^{\infty}f'(x)e^{-iux}\,dx = \bigl[f(x)e^{-iux}\bigr]_{-\infty}^{\infty}+iu\int_{-\infty}^{\infty}f(x)e^{-iux}\,dx = iuF(u)$\\[2pt]
  {\footnotesize\textcolor{gray}{$|x|\to\infty$ で $f(x)\to0$ を仮定。繰り返し適用で $\mathcal{F}[f^{(n)}]=(iu)^nF(u)$}}
  \\[6pt]\textcolor{red}{\textbf{答}}：証明終
\end{frame}

\begin{frame}{問5　}
  \begin{exampleblock}{問題}
    性質 III および VI を証明せよ
  \end{exampleblock}
\end{frame}

\begin{frame}{問5　解答}
  $III:\; \mathcal{F}[e^{iax}f(x)] = \displaystyle\int_{-\infty}^{\infty} e^{iax}f(x)e^{-iux}\,dx = \int_{-\infty}^{\infty} f(x)e^{-i(u-a)x}\,dx = F(u-a)$\\[8pt]
  $VI:\; F^{(n)}(u) = \dfrac{d^{n}}{du^{n}}\displaystyle\int_{-\infty}^{\infty} f(x)e^{-iux}\,dx = \int_{-\infty}^{\infty} f(x)(-ix)^{n}e^{-iux}\,dx = \mathcal{F}[(-ix)^{n}f(x)]$
  \\[8pt]\textcolor{red}{\textbf{答}}：証明終
\end{frame}

\begin{frame}{たたみこみ}
  \begin{block}{定義}
  \[
    (f * g)(x) = \int_{-\infty}^{\infty} f(t)\,g(x-t)\,dt
  \]
  \end{block}
  {\footnotesize\textcolor{gray}{$f * g$ を $f$ と $g$ のたたみこみ（合成積）という}}
\end{frame}

\begin{frame}{たたみこみのフーリエ変換}
  \begin{block}{定理}
  \[
    \mathcal{F}[f * g] = F(u)\cdot G(u)
  \]
  \end{block}
  {\footnotesize\textcolor{gray}{たたみこみのフーリエ変換は、それぞれのフーリエ変換の積に等しい}}
\end{frame}

\begin{frame}{問6　}
  \begin{exampleblock}{問題}
    $\mathcal{F}[f(x)] = F(u)$，$\mathcal{F}[g(x)] = G(u)$ のとき，次の等式を証明せよ\\[6pt]
    $\mathcal{F}[f(x)g(x)] = \dfrac{1}{2\pi} F(u) * G(u)$
  \end{exampleblock}
\end{frame}

\begin{frame}{問6　解答}
  $g(x) = \dfrac{1}{2\pi}\displaystyle\int_{-\infty}^{\infty} G(v)\,e^{ivx}\,dv$ を $\mathcal{F}[f(x)g(x)]$ に代入\\[6pt]
  $= \dfrac{1}{2\pi}\displaystyle\int_{-\infty}^{\infty} G(v)\!\left[\int_{-\infty}^{\infty} f(x)e^{-i(u-v)x}\,dx\right]dv = \dfrac{1}{2\pi}\int_{-\infty}^{\infty} G(v)F(u-v)\,dv = \dfrac{1}{2\pi}F*G$
  \\[8pt]\textcolor{red}{\textbf{答}}：証明終
\end{frame}

\begin{frame}{問7　}
  \begin{exampleblock}{問題}
    96ページの性質を用いて，次の関数のフーリエ変換を求めよ\\[6pt]
    $(1)\; xe^{-x^{2}/2}$\\[4pt]
    $(2)\; x^{2}e^{-x^{2}/2}$
  \end{exampleblock}
\end{frame}

\begin{frame}{問7(1)　解答}
  性質 $VI:\; \mathcal{F}[(-ix)^{n}f(x)] = F^{(n)}(u)$，$\mathcal{F}[e^{-x^{2}/2}] = \sqrt{2\pi}\,e^{-u^{2}/2}$\\[6pt]
  $n=1:\; \mathcal{F}[(-ix)e^{-x^{2}/2}] = \dfrac{d}{du}\!\left[\sqrt{2\pi}\,e^{-u^{2}/2}\right] = -\sqrt{2\pi}\,u\,e^{-u^{2}/2}$\\[6pt]
  $\mathcal{F}[xe^{-x^{2}/2}] = \dfrac{1}{-i}\cdot(-\sqrt{2\pi}\,u\,e^{-u^{2}/2}) = -i\sqrt{2\pi}\,u\,e^{-u^{2}/2}$
  \\[8pt]\textcolor{red}{\textbf{答}}：$\mathcal{F}[xe^{-x^{2}/2}] = -i\sqrt{2\pi}\,u\,e^{-u^{2}/2}$
\end{frame}

\begin{frame}{問7(2)　解答}
  $n=2:\; (-ix)^{2} = -x^{2}$\\[6pt]
  $\mathcal{F}[-x^{2}e^{-x^{2}/2}] = \dfrac{d^{2}}{du^{2}}\!\left[\sqrt{2\pi}\,e^{-u^{2}/2}\right] = \sqrt{2\pi}(u^{2}-1)e^{-u^{2}/2}$\\[6pt]
  $\mathcal{F}[x^{2}e^{-x^{2}/2}] = \sqrt{2\pi}(1-u^{2})\,e^{-u^{2}/2}$
  \\[8pt]\textcolor{red}{\textbf{答}}：$\mathcal{F}[x^{2}e^{-x^{2}/2}] = \sqrt{2\pi}(1-u^{2})\,e^{-u^{2}/2}$
\end{frame}

\begin{frame}{問8　}
  \begin{exampleblock}{問題}
    次の等式を証明せよ（$b$ は正の定数）\\[6pt]
    $\mathcal{F}^{-1}[e^{-bu^{2}}] = \dfrac{1}{2\sqrt{\pi b}}\,e^{-x^{2}/(4b)}$
  \end{exampleblock}
\end{frame}

\begin{frame}{問8　解答}
  p.98 公式 $\mathcal{F}[e^{-ax^{2}}] = \sqrt{\dfrac{\pi}{a}}\,e^{-u^{2}/(4a)}$ で $b = \dfrac{1}{4a}$ とおくと $a = \dfrac{1}{4b}$\\[6pt]
  $\mathcal{F}[e^{-x^{2}/(4b)}] = \sqrt{4\pi b}\,e^{-bu^{2}} = 2\sqrt{\pi b}\,e^{-bu^{2}}$\\[6pt]
  両辺を $2\sqrt{\pi b}$ で割って逆変換：$\mathcal{F}^{-1}[e^{-bu^{2}}] = \dfrac{e^{-x^{2}/(4b)}}{2\sqrt{\pi b}}$
  \\[8pt]\textcolor{red}{\textbf{答}}：証明終
\end{frame}

\begin{frame}{例題4　（p.\,99）}
  \begin{exampleblock}{問題}
    次の等式を証明せよ\\[4pt]
    $(1)\; \mathcal{F}[e^{-a^{2}x^{2}} * e^{-b^{2}x^{2}}] = \dfrac{\pi}{ab}\,e^{-u^{2}(a^{2}+b^{2})/(4a^{2}b^{2})}$\\[6pt]
    $(2)\; e^{-a^{2}x^{2}} * e^{-b^{2}x^{2}}$ を求めよ
  \end{exampleblock}
\end{frame}

\begin{frame}{例題4　解答}
  $(1)$ p.\,98 公式 $\mathcal{F}[e^{-cx^{2}}] = \sqrt{\pi/c}\,e^{-u^{2}/(4c)}$ より\\[4pt]
  $\mathcal{F}[e^{-a^{2}x^{2}}] = \dfrac{\sqrt{\pi}}{a}e^{-u^{2}/(4a^{2})},\quad \mathcal{F}[e^{-b^{2}x^{2}}] = \dfrac{\sqrt{\pi}}{b}e^{-u^{2}/(4b^{2})}$\\[4pt]
  たたみこみ定理を適用：$\mathcal{F}[e^{-a^{2}x^{2}}*e^{-b^{2}x^{2}}] = \dfrac{\pi}{ab}\,e^{-u^{2}(a^{2}+b^{2})/(4a^{2}b^{2})}$\\[6pt]
  $(2)$ 結果を $\sqrt{\pi/c}\,e^{-u^{2}/(4c)}$ の形と比較：$c = \dfrac{a^{2}b^{2}}{a^{2}+b^{2}}$\\[4pt]
  $e^{-a^{2}x^{2}}*e^{-b^{2}x^{2}} = \dfrac{\sqrt{\pi}}{\sqrt{a^{2}+b^{2}}}\,e^{-(a^{2}b^{2}/(a^{2}+b^{2}))x^{2}}$
  \\[6pt]\textcolor{red}{\textbf{答}}：証明終
\end{frame}

\begin{frame}{問9　}
  \begin{exampleblock}{問題}
    次の等式を証明せよ（$*$ はたたみこみ）\\[6pt]
    $(1)\; \mathcal{F}[e^{-x^{2}/2} * xe^{-x^{2}/2}] = -2\pi iu\,e^{-u^{2}}$\\[4pt]
    $(2)\; e^{-x^{2}/2} * xe^{-x^{2}/2} = \dfrac{\sqrt{\pi}}{2}\,xe^{-x^{2}/4}$
  \end{exampleblock}
\end{frame}

\begin{frame}{問9(1)　解答}
  たたみこみ定理：$\mathcal{F}[f*g] = F(u)\cdot G(u)$\\[6pt]
  $\mathcal{F}[e^{-x^{2}/2}] = \sqrt{2\pi}\,e^{-u^{2}/2},\quad \mathcal{F}[xe^{-x^{2}/2}] = -i\sqrt{2\pi}\,u\,e^{-u^{2}/2}$（問7(1)）\\[6pt]
  積：$\sqrt{2\pi}\,e^{-u^{2}/2}\cdot(-i\sqrt{2\pi}\,u\,e^{-u^{2}/2}) = -2\pi iu\,e^{-u^{2}}$
  \\[8pt]\textcolor{red}{\textbf{答}}：$\mathcal{F}[e^{-x^{2}/2}*xe^{-x^{2}/2}] = -2\pi iu\,e^{-u^{2}}$
\end{frame}

\begin{frame}{問9(2)　解答}
  $\mathcal{F}^{-1}[-2\pi iu\,e^{-u^{2}}]$ を計算\\[6pt]
  問8（$b=1$）より $\mathcal{F}^{-1}[e^{-u^{2}}] = \dfrac{e^{-x^{2}/4}}{2\sqrt{\pi}}$\\[6pt]
  $I(x)=\displaystyle\int_{-\infty}^{\infty}e^{-u^{2}}e^{iux}\,du = \sqrt{\pi}\,e^{-x^{2}/4}$ を微分：$\int_{-\infty}^{\infty}ue^{-u^{2}}e^{iux}\,du = \dfrac{ix}{2}\sqrt{\pi}\,e^{-x^{2}/4}$\\[6pt]
  $\mathcal{F}^{-1}[-2\pi iu\,e^{-u^{2}}] = (-i)\cdot\dfrac{1}{2\pi}\cdot2\pi\cdot\dfrac{ix}{2}\sqrt{\pi}\,e^{-x^{2}/4} = \dfrac{\sqrt{\pi}}{2}\,x\,e^{-x^{2}/4}$
  \\[8pt]\textcolor{red}{\textbf{答}}：証明終
\end{frame}

\end{document}
